% !TeX program = lualatex
% Compile twice: lualatex 181.tex
% All references and prose are embedded; no BibTeX database or external inputs.
\documentclass[11pt,a4paper]{article}
\usepackage[margin=24mm,headheight=14pt]{geometry}
\usepackage{fontspec}
\setmainfont{Latin Modern Roman}
\setsansfont{Latin Modern Sans}
\setmonofont{Latin Modern Mono}
\usepackage{amsmath,amssymb,mathtools}
\usepackage{unicode-math}
\setmathfont{Latin Modern Math}
\usepackage{microtype}
\usepackage{enumitem,xurl,needspace}
\usepackage[dvipsnames]{xcolor}
\usepackage{hyperref}
\hypersetup{unicode=true,colorlinks=true,linkcolor=MidnightBlue,urlcolor=MidnightBlue,
 pdftitle={500 Mathematical Problem Statements},pdfauthor={Source-annotated compilation},bookmarksnumbered=true}
\usepackage{bookmark}
\usepackage{tocloft}
\setlength{\cftsecnumwidth}{3em}
\cftsetpnumwidth{2.5em}
\cftsetrmarg{4em}
\usepackage{titlesec}
\titleformat{\section}{\Large\bfseries\raggedright}{\thesection}{0.7em}{}
\usepackage{fancyhdr}
\pagestyle{fancy}\fancyhf{}
\fancyhead[L]{\small\sffamily Mathematical problem statements}
\fancyhead[R]{\small Source edition 22}
\fancyfoot[C]{\thepage}
\setlength{\parindent}{0pt}
\setlength{\parskip}{5pt plus 1pt}
\setlength{\emergencystretch}{3em}
\setcounter{tocdepth}{1}
\setcounter{secnumdepth}{1}
\setlist[itemize]{leftmargin=3em,itemsep=5pt,topsep=4pt,parsep=0pt}
\allowdisplaybreaks
\newcommand{\N}{\mathbb{N}}
\newcommand{\Z}{\mathbb{Z}}
\newcommand{\Q}{\mathbb{Q}}
\newcommand{\R}{\mathbb{R}}
\newcommand{\C}{\mathbb{C}}
\newcommand{\F}{\mathbb{F}}
\newcommand{\A}{\mathbb{A}}
\newcommand{\PP}{\mathbb P}
\newcommand{\E}{\mathbb{E}}
\newcommand{\eps}{\varepsilon}
\newcommand{\dd}{\,\mathrm d}
\newcommand{\sref}[2]{\hyperlink{source:#1:#2}{[S#2]}}
\newcommand{\eref}[2]{\hyperlink{extra:#1:#2}{[E#2]}}
% Turn the next switch off for a shorter reading copy. The quotations preserve
% every exactTarget field, including qualifications omitted from a short summary.
\newif\ifshowcatalogue
\showcataloguetrue
\newcommand{\cataloguescope}[1]{\ifshowcatalogue
 \paragraph{Exact scope in the supplied catalogue.}
 \begin{quote}\small #1\end{quote}\fi}

% Standalone notebook wrapper; original problem section retained verbatim.
\numberwithin{equation}{section}
\hypersetup{pdftitle={Research attempt -- problem 181},
 pdfauthor={Research notebook; source compilation retained}}
\fancyhead[L]{\small\sffamily Research notebook}
\fancyhead[R]{\small\sffamily Problem 181 | 27 September 2026}
\begin{document}
\setcounter{section}{180}
% ================================================================
% problemId: problem.l-space-conjecture
% source releaseRank: 181; source releaseStatus: open_with_solved_subcases
% exactTarget SHA-256: dda5d69cbb0816417849e6b26e8c9c861e6a69f69cd768eb1a9f44b8c491c7c9
\clearpage
\section{\texorpdfstring{L-space Conjecture}{L-space Conjecture}}\label{181}
\subsection{Definitions and mathematical statement}
A rational homology three-sphere $Y$ has $H_*(Y,\Q)\cong H_*(S^3,\Q)$; irreducibility means every embedded two-sphere bounds a ball. An L-space has
$\operatorname{rank}\widehat{HF}(Y)=|H_1(Y,\Z)|$ in the Heegaard Floer convention of the source. A group is left-orderable if it has a strict total order preserved by left multiplication. A codimension-one foliation is coorientable if its normal line bundle is orientable, and taut if every leaf meets a closed transversal. For closed oriented irreducible rational homology three-spheres, the conjecture equates: not being an L-space; having a nontrivial left-orderable fundamental group; and admitting a coorientable taut foliation. The explicit nontriviality clause excludes the trivial-group convention from changing the statement.
\par\smallskip\textit{Formulation and concept references:} \sref{181}{1}.
\cataloguescope{Prove or disprove: for a closed, oriented, irreducible rational homology 3-sphere Y, the following are equivalent: (1) Y is not a Heegaard Floer L-space; (2) pi\_1(Y) is nontrivial and admits a left-invariant strict total order; (3) Y admits a coorientable taut foliation.}
\subsection{Short English statement}
Are three descriptions of a rational homology three-sphere equivalent: nonminimal Floer homology, an orderable nontrivial fundamental group, and a coorientable taut foliation?
% BEGIN RESEARCH ADDITION ATTEMPT-181
\subsection{Research attempt}

\paragraph{Current frontier.}
The three conditions in the original entry must remain separate. The
implication from a coorientable taut foliation to nonminimal Heegaard Floer
homology and the graph-manifold results are reviewed in \sref{181}{1}.
A 2026 manuscript of Boyer--Gordon--Hu proves circular orderability for
closed orientable irreducible toroidal three-manifolds and constructs
finite cyclic covers with left-orderable fundamental groups
\eref{181}{1}, Theorems 1.1 and 1.4. Neither conclusion is left-orderability
of the original group. Their 2025 relative-conjecture manuscript was
withdrawn on 18 August 2025 because of an error in the proof of Theorem 5.7;
no conclusion from that withdrawn argument is used \eref{181}{2}.
These source distinctions motivate the lifting obstruction examined below.

\paragraph{Attack route.}
A Floer-theoretic route needs to turn a nonminimal Floer group into a
geometric or dynamical object, rather than merely compute its rank. A
foliation route needs enough control of its transverse dynamics to produce
an order. A direct group route could attempt to exclude every finite
obstruction to an order. We combine the latter two: first make the finite
obstruction criterion explicit, then identify the exact integral class
that must vanish in a circle-action lifting argument. Rational vanishing
or passage to a finite cover is tested rather than presumed sufficient.

\paragraph{Attempt: left-orderability has finite sign obstructions.}
Let $G$ be a countable group with identity $e$. The following precise
criterion is useful:
\begin{equation}\label{eq181-signs}
 G\text{ is left-orderable}\quad\Longleftrightarrow\quad
 \forall F\subset G\setminus\{e\}\text{ finite},\ \exists
 (\eps_g)_{g\in F}\in\{\pm1\}^{F}:
 e\notin\langle g^{\eps_g}:g\in F\rangle_+.
\end{equation}
The subscript $+$ means nonempty words using only the indicated elements;
it is a semigroup, not the subgroup they generate.

\textit{Proof.} A left order selects the positive one of $g,g^{-1}$,
and products of positive elements stay positive, so the right side is
necessary. For sufficiency enumerate the nonidentity elements
$g_1,g_2,\ldots$. At level $r$ of a binary tree retain the sign choices
whose generated positive semigroup does not contain $e$. A retained node
has a retained prefix. The hypothesis gives at least one node at every
level, and finite branching gives an infinite branch. Let $P$ be the
semigroup generated by the selected signed elements along this branch.
A word witnessing $e\in P$ would already occur at a finite level, so
$e\notin P$. For every $g\ne e$, either $g$ or $g^{-1}$ belongs to $P$;
not both, since their product would be $e$. Thus
\[
               G=P\sqcup\{e\}\sqcup P^{-1},\qquad PP\subseteq P.
\]
Define $g<h$ when $g^{-1}h\in P$. Trichotomy and transitivity follow from
these two properties, and $(ag)^{-1}(ah)=g^{-1}h$ proves left invariance.
The finite-group case is obtained from the same finite-prefix argument,
or by enumerating with repetitions. \hfill$\square$

Consequently, if $G$ is not left-orderable, some finite list of nonidentity
elements has the following certificate: for each of its finitely many sign
choices there is a positive word equal to $e$. The words may have very
different lengths. The proof gives no computable uniform length bound and
no efficient procedure to find them. In applying such a certificate one
must also justify that the listed elements really are nonidentity in $G$.
A torsion element of order $m>1$ is the simplest example: its $m$th power,
with either sign, gives the obstruction.

Thus a possible topological route to $(1)\Rightarrow(2)$ would be a theorem
excluding all these finite obstructions for fundamental groups of non-L-spaces.
Nothing in \eqref{eq181-signs} supplies that connection to Floer homology.
Merely checking many finite presentations or many sign lists does not
establish the universal absence of an obstruction.

\paragraph{Attempt: lift a faithful circle action, with its full obstruction.}
Suppose a faithful representation
$\rho:G\to\operatorname{Homeo}_+(S^1)$ has been constructed. Choose lifts
$\widetilde\rho(g)$ to increasing homeomorphisms of $\R$ commuting with
$T(x)=x+1$. The failure of these lifts to be a homomorphism is an integer
cocycle:
\[
 \widetilde\rho(g)\widetilde\rho(h)
       =T^{c(g,h)}\widetilde\rho(gh),\qquad c(g,h)\in\Z.
\]
Associativity gives the cocycle identity. Changing the lifts by integer
translations changes $c$ by an integer coboundary. Its class
$e(\rho)\in H^2(G;\Z)$ is the integral Euler class of the action.

\textbf{Lifting lemma.} If $e(\rho)=0$, then $G$ is left-orderable.

\textit{Proof.} Write $c$ as a coboundary and change each lift by the
corresponding integer translation. The modified lifts form a homomorphism
$\widehat\rho:G\to\operatorname{Homeo}_+(\R)$. It is faithful, since an
element acting trivially upstairs acts trivially downstairs, where $\rho$
is faithful. To construct an order explicitly, enumerate a dense set
$q_1,q_2,\ldots$ in $\R$. For distinct $g,h$, continuity and faithfulness
ensure that some $q_i$ distinguishes their actions. Compare
$\widehat\rho(g)(q_i)$ and $\widehat\rho(h)(q_i)$ at the first such index.
This lexicographic comparison is a strict total order, and composition on
the left by an increasing homeomorphism preserves the comparison and all
preceding equalities. Hence it is left invariant. \hfill$\square$

This proof does not require a free action on $\R$, or that every group
element move every point. It requires a faithful action and the vanishing
of an \emph{integral} cocycle. Similar Euler-class lifting mechanisms
occur in the actual three-manifold arguments of \eref{181}{1}, for example
its proof of Lemma 4.4; the general hypotheses above are isolated here
instead of assumed for all manifolds in the target.

\paragraph{Stress test: rational vanishing is not enough.}
A rational homology sphere has $H^2(Y;\Q)=0$, but its integral
$H^2(Y;\Z)$ can contain torsion. Even identifying this with group
cohomology requires an asphericity hypothesis; it is not true for every
space solely because its fundamental group is specified. For an
aspherical $Y$, integral group cohomology equals that of $Y$, but rational
vanishing still leaves exactly the torsion obstruction visible above.

The faithful action of $C_m$ on $S^1$ by rotations through multiples of
$1/m$ is a direct test. Its Euler class vanishes rationally, but it cannot
lift to a faithful action on $\R$: an increasing homeomorphism of finite
order is the identity. Indeed, if $h(x)>x$, its successive iterates are
strictly increasing, and if $h(x)<x$, they are strictly decreasing; neither
is compatible with finite order. This also proves that $C_m$ is not
left-orderable. Thus rational Euler-class vanishing cannot replace the
hypothesis of the lifting lemma.

Likewise, an orderable finite-index subgroup need not make the whole group
orderable. The trivial subgroup has finite index in $C_m$, while the whole
group fails the preceding test. This example does not purport to capture
all extra geometry of an irreducible toroidal manifold, but it disproves
the bare group-theoretic implication that would otherwise incorrectly
upgrade the cyclic-cover conclusion of \eref{181}{1}.

\paragraph{Attempt: combine the two diagnostics.}
A proposed construction from a taut foliation could produce circle dynamics
and then try to lift. The exact sufficient input is now explicit: a
\emph{faithful} circle action with integrally trivial Euler class, or some
other faithful real-line action. When such an input is available, the
lifting proof yields a positive cone and therefore excludes every finite
certificate in \eqref{eq181-signs}. A circle action with a nontrivial kernel
instead orders at most the image by this method; an additional argument
would have to handle the kernel compatibly. No such kernel or Euler-class
control has been established here for an arbitrary coorientable taut
foliation in the catalogue's class.

Conversely, trying to use an arbitrary left order to produce a taut
foliation requires a separate realization theorem in the topology of $Y$.
The ordered group by itself contains no constructed plane field,
integrability condition, or closed transversals. The criterion
\eqref{eq181-signs} therefore does not silently establish
$(2)\Rightarrow(3)$. Nor does the known foliation-to-Floer implication
close a cycle of implications when its other arrows remain missing.

The explicit nontriviality convention is also respected: the trivial group
admits a formal left order, but condition (2) of this entry rules it out.
Nothing above changes that convention for $S^3$. Infinite fixed sets of
an action or a finite cover's foliation are not being substituted for a
foliation on the original manifold.

\paragraph{Outcome.}
\textbf{Outcome: Exploratory approach with unresolved gap.}

The attempt proves the finite sign-obstruction criterion and a precise
Euler-class lifting lemma, then exhibits torsion examples defeating two
natural shortcuts. These are group-theoretic tools and diagnostics, not
new solved three-manifold cases or claims of priority for these standard
constructions.

The missing geometric input is a mechanism tying nonminimal Floer homology
or a taut foliation to the required faithful, integrally liftable dynamics,
together with the converse realization where needed. The current
circular-order and finite-cover results do not supply those assertions for
the original manifold. No full equivalence or counterexample is obtained.
% END RESEARCH ADDITION ATTEMPT-181
\subsection{Sources}
\begin{itemize}\raggedright
\item[\textnormal{[S1]}]\hypertarget{source:181:1}{} J. Hanselman, "L-spaces, foliations, and left-orderability," Princeton RTG lecture, notes by J. Van Dyke; conjecture due to Boyer-Gordon-Watson (Math. Ann. 356 (2013)) and Juhasz.
\url{https://web.ma.utexas.edu/users/vandyke/notes/princeton_rtg_notes/hanselman/hanselman.pdf}.
{\small\textit{Catalogue formulation source.}}
% BEGIN RESEARCH ADDITION SOURCES-181
\item[\textnormal{[E1]}]\hypertarget{extra:181:1}{} Steven Boyer, Cameron
McA. Gordon and Ying Hu, \emph{Toroidal 3-manifolds have circularly orderable
fundamental groups}, arXiv:2607.10529v1, 12 July 2026, Theorems 1.1 and 1.4,
and Lemma 4.4. \url{https://arxiv.org/html/2607.10529v1}.
{\small\textit{Consulted 27 September 2026. Recent manuscript; circular
order and left-orderable finite covers are distinguished from an order
on the original group.}}
\item[\textnormal{[E2]}]\hypertarget{extra:181:2}{} Steven Boyer, Cameron
McA. Gordon and Ying Hu, \emph{Slope detection, taut foliations, and the
relative L-space conjecture}, arXiv:2508.06395v2, 18 August 2025,
\textbf{withdrawn}. \url{https://arxiv.org/abs/2508.06395v2}.
{\small\textit{Status only, checked 27 September 2026. The withdrawal
notice identifies an error in the proof of Theorem 5.7. No theorem from
that manuscript is used.}}
% END RESEARCH ADDITION SOURCES-181
\end{itemize}

\end{document}
